\chapter{Statistical reference laws and uncertainty calculations}
\label{sec:shared-statistical-methods}
\label{sec:simulation-statistical-methods}

These results concern ideal independent observations from the stated model.
They justify reference laws and interval calculations, rather than proving
that finite-state pseudorandom output is independent or that the model fits
real tests. The input theorems are Gaussian affine closure and the Gaussian
characteristic function \citep[Theorems 3.3.1--3.3.2]{gallager2013},
orthonormal-basis extension \citep[Theorem 6.36]{axler2024}, and linearity,
independence and elementary probability inequalities
\citep[Sections 1.6, 2.1 and 3.4]{durrett2019}. The derivations below do
not depend on a successful software test.

\begin{proposition}[Normal sample moments and the chi-square pivot]
\label{sm:normal-sample}
For an integer \(m\geq2\), independent \(Y_i\sim N(\mu,v)\), \(v>0\),
the sample mean is \(N(\mu,v/m)\). With
\(S^2=\sum_i(Y_i-\bar Y)^2/(m-1)\),
\((m-1)S^2/v\sim\chi^2_{m-1}\), independently of \(\bar Y\), and
\(\E S^2=v\).
\end{proposition}
\begin{proof}
Finite linearity and zero cross-covariances give the mean and variance;
Gaussian affine closure gives the law. Set \(Z_i=(Y_i-\mu)/\sqrt v\),
and extend the unit vector \(m^{-1/2}(1,\ldots,1)\) to an orthonormal
basis. Let \(Q\) have this vector as its first row. The transformed
Gaussian \(T=QZ\) has covariance \(QQ^{\mathsf T}=I\); its characteristic
function factors, so all its coordinates are independent standard normals.
Its first coordinate is \(\sqrt m(\bar Y-\mu)/\sqrt v\). Orthogonality
preserves squared norm, hence
\[
 \frac{1}{v}\sum_i(Y_i-\bar Y)^2
 =\sum_iZ_i^2-T_1^2=\sum_{j=2}^mT_j^2.
\]
The right side is chi-square by definition, independent of \(T_1\).
Its expectation is \(m-1\), proving unbiasedness of \(S^2\).
This argument supplies the normal and chi-square laws directly from the
referenced Gaussian and linear-algebra theorems.
\end{proof}

\begin{corollary}[SD intervals and fixed-parameter reference bands]
\label{sm:sd-interval}
Let \(q_u\) denote a \(\chi^2_{m-1}\) quantile and \(0<\alpha<1\).
An exact \(1-\alpha\) confidence interval for the unknown SD under this
normal model is
\[
 \left[\sqrt{\frac{(m-1)S^2}{q_{1-\alpha/2}}},
       \sqrt{\frac{(m-1)S^2}{q_{\alpha/2}}}\right].
\]
For a fixed theoretical \(v\), a reference interval for \(S^2\) is instead
\([vq_{\alpha/2}/(m-1),vq_{1-\alpha/2}/(m-1)]\). The reference
mean discrepancy in standard-error units is
\((\bar Y-\mu)/\sqrt{v/m}\sim N(0,1)\).
\end{corollary}
\begin{proof}
The continuous chi-square law gives
\(P(q_{\alpha/2}\leq(m-1)S^2/v\leq q_{1-\alpha/2})=1-\alpha\).
Solving these inequalities for \(v\), then taking the increasing positive
square root, gives the SD interval. Solving for \(S^2\), keeping \(v\)
fixed, gives the reference interval. The last assertion is the standardized
sample-mean law in Proposition~\ref{sm:normal-sample}. If \(v=0\), every
observation equals \(\mu\) almost surely, so this division-based pivot is
inapplicable; constant output is checked directly. SciPy evaluates the
quantiles \citep{scipystats}; the algebra, rather than the library call,
establishes the probability statement.
\end{proof}

\begin{proposition}[Bernoulli counts and sampling precision]
\label{sm:binomial}
For an integer \(m\geq1\) and \(m\) independent event indicators of common
probability \(p\in[0,1]\), the
count \(K\) is binomial, \(\E K=mp\), and
\(\Var(K/m)=p(1-p)/m\).
\end{proposition}
\begin{proof}
Independence makes the probability generating function of their sum
\((1-p+pt)^m\). The binomial theorem gives
\(P(K=j)=\binom mjp^j(1-p)^{m-j}\). Each indicator has mean \(p\)
and variance \(p(1-p)\), because its square equals itself. Summing means
and independent variances, then dividing by \(m^2\), proves the claims.
These are finite-product independence and linearity
\citep[Sections 1.6 and 2.1]{durrett2019}.
\end{proof}

\begin{proposition}[Wilson score endpoints and their coverage scope]
\label{sm:wilson}
For an integer \(m\geq1\), an observed count \(K\in\{0,\ldots,m\}\),
\(0<\alpha<1\), \(\hat p=K/m\) and
\(z=\Phi^{-1}(1-\alpha/2)\), inversion
of the normal score inequality gives endpoints
\[
 \frac{\hat p+z^2/(2m)\ \mathord\pm\
 z\sqrt{\hat p(1-\hat p)/m+z^2/(4m^2)}}{1+z^2/m}.
\]
The endpoint algebra is exact; nominal \(1-\alpha\) coverage uses a normal
score approximation and is not an exact finite-sample coverage theorem.
\end{proposition}
\begin{proof}
Squaring \(|\hat p-p|\leq z\sqrt{p(1-p)/m}\) gives
\[
 (m+z^2)p^2-(2m\hat p+z^2)p+m\hat p^2\leq0.
\]
The positive leading coefficient means precisely the interval between its
two real roots. Its discriminant is
\(4mz^2\hat p(1-\hat p)+z^4\). The quadratic formula, dividing
numerator and denominator by \(2m\), gives the stated endpoints. For
\(0<p<1\), the classical independent-identically-distributed central
limit theorem implies
\(\sqrt m(\hat p-p)/\sqrt{p(1-p)}\Rightarrow N(0,1)\), giving limiting
coverage at its continuity points, not equality for fixed \(m\).
This is the score construction of \citet{wilson1927}; the limit theorem is
\citet[Theorem 3.4.1]{durrett2019}. The implementation reuses
SciPy's binomial-result interval routine \citep{scipyBinomtest}.
\end{proof}

\begin{proposition}[Exact equal-tail binomial intervals]
\label{sm:clopper-pearson}
For an integer \(m\geq1\), a count \(j\in\{0,\ldots,m\}\), and
\(0<\alpha<1\), define \(L(0)=0\), \(U(m)=1\); for \(j>0\),
let \(L(j)\) solve \(P_{L(j)}(K\geq j)=\alpha/2\), and for \(j<m\),
let \(U(j)\) solve \(P_{U(j)}(K\leq j)=\alpha/2\). Then
\(P_p(L(K)\leq p\leq U(K))\geq1-\alpha\) for every \(p\in[0,1]\).
\end{proposition}
\begin{proof}
For \(0<j\leq m\), differentiating the finite binomial sum and
telescoping gives
\[
 \frac{d}{dp}P_p(K\geq j)
 =m\binom{m-1}{j-1}p^{j-1}(1-p)^{m-j}>0\quad(0<p<1).
\]
The tail runs continuously from zero to one, so \(L(j)\) exists uniquely;
the lower tail is strictly decreasing and gives \(U(j)\) similarly.
For \(0<j<m\), at \(p=L(j)\) the lower tail is at least
\(1-\alpha/2>\alpha/2\); its decrease implies \(U(j)>L(j)\).
The endpoint counts are ordered directly, so these really are intervals.
For a fixed \(p\), \(L(j)>p\) implies \(P_p(K\geq j)<\alpha/2\).
Such \(j\)'s form an upper terminal set, because upper tails decrease
with \(j\). If nonempty its probability is the tail at its smallest
member and is at most \(\alpha/2\). The analogous lower terminal set
\(U(j)<p\) has probability at most \(\alpha/2\). The union bound makes
total noncoverage at most \(\alpha\); endpoint parameters follow directly.
Discreteness explains why the inequality can be strict. This is the
test-inversion construction of \citet{clopper1934}, with its coverage
proof supplied here. Equivalently the endpoints are beta quantiles:
the derivative above is the \(\mathrm{Beta}(j,m-j+1)\) density and its
tail has value zero at \(p=0\), so integration gives the same cumulative
distribution; the complementary tail gives the upper endpoint.
\end{proof}

\begin{proposition}[Conditional frequencies with random denominators]
\label{sm:conditional-count}
For independent repetitions of a four-cell event experiment, condition on
\(N_A=\ell>0\) observations in event \(A\), where \(P(A)>0\).
The count also in \(B\) is then \(\mathrm{Binomial}(\ell,P(B\mid A))\).
Thus a binomial interval conditional on its observed denominator has its
stated exact or approximate coverage scope. A zero denominator gives no
defined conditional-frequency estimate.
\end{proposition}
\begin{proof}
Let \(q=P(A)\) and \(r=P(B\mid A)\). Grouping cells into
\(A\cap B,A\setminus B,A^c\), their probabilities are
\(qr,q(1-r),1-q\). Product independence and the multinomial counting
coefficient give
\[
 P(N_{AB}=j,N_A=\ell)
 =\binom m\ell\binom\ell j q^\ell(1-q)^{m-\ell}
           r^j(1-r)^{\ell-j}.
\]
Divide by \(P(N_A=\ell)=\binom m\ell q^\ell(1-q)^{m-\ell}\)
when positive. The remaining factor is the stated binomial mass. Cases
with \(q=1\) follow directly with \(\ell=m\). Conditional coverage for
each positive \(\ell\) also gives coverage averaged over positive
denominators by total probability. No convention assigning a number to
\(0/0\) supplies data when \(\ell=0\).
\end{proof}

\begin{proposition}[Pooling counts and moments]
\label{sm:pooling}
For a nonempty finite collection of groups of integer sizes \(m_j\geq2\),
means \(\bar Y_j\) and sample variances
\(S_j^2\), let \(M=\sum_jm_j\) and
\(\bar Y=M^{-1}\sum_jm_j\bar Y_j\). The pooled sample variance satisfies
\[
 (M-1)S^2=\sum_j(m_j-1)S_j^2+\sum_jm_j(\bar Y_j-\bar Y)^2.
\]
Independent equal-design Bernoulli groups pool to a binomial count of size
\(M\); independence is required for that probability law, not for the
displayed algebraic identity.
\end{proposition}
\begin{proof}
Within group \(j\), write
\(Y_{ji}-\bar Y=(Y_{ji}-\bar Y_j)+(\bar Y_j-\bar Y)\).
Squaring and summing, its cross term vanishes because
\(\sum_i(Y_{ji}-\bar Y_j)=0\). Summing groups gives the identity.
For counts, independence multiplies their generating functions to
\(\prod_j(1-p+pt)^{m_j}=(1-p+pt)^M\), proving the binomial law by
Proposition~\ref{sm:binomial}. Sharing random numbers across cases does not
establish this independence premise; in the overlapping-factor designs used
here the rows cannot be counted as additional independent replications.
Averaging group SDs generally does
not satisfy the displayed within-plus-between identity.
\end{proof}

\begin{proposition}[Simultaneous diagnostics and rare-event precision]
\label{sm:multiple}
If \(J\) fixed diagnostics have false-alarm probabilities at most
\(\alpha_j\), the probability of any false alarm is at most
\(\sum_{j=1}^J\alpha_j\), regardless of their dependence.
\end{proposition}
\begin{proof}
The pointwise bound \(1_{\cup_j F_j}\leq\sum_j1_{F_j}\), integrated,
is the union bound. Taking each \(\alpha_j=\alpha/J\) therefore gives
a family bound \(\alpha\). The same argument applies to two tails per
row, counting both in \(J\). Finite diagnostics remain observations
under a reference model and cannot establish an entire distribution.
The bound is the elementary subadditivity theorem
\citep[Section 1.1]{durrett2019}.
\end{proof}

\begin{corollary}[Zero counts, relative precision and sample-size scaling]
\label{sm:rare}
For independent Bernoulli trials, \(P(K=0)=(1-p)^m\). For \(p>0\) the
relative standard error of \(K/m\) is \(\sqrt{(1-p)/(mp)}\); requiring
it at most \(\epsilon>0\) means \(m\geq(1-p)/(p\epsilon^2)\).
At fixed \(p\in(0,1)\), multiplying \(m\) by four halves the standard
error. Zero successes do not prove \(p=0\).
\end{corollary}
\begin{proof}
The first statement is the binomial mass at zero. Divide the SD in
Proposition~\ref{sm:binomial} by \(p\), square the nonnegative inequality,
and rearrange to get the sample-size bound. The square-root dependence
proves the scaling. Every \(0<p<1\) assigns positive probability to
zero successes for finite \(m\), proving the final statement. A desired
95\% zero-count upper confidence bound, for example, follows from
\((1-U)^m=.025\) for the equal-tail construction, not from setting
\(p\) equal to zero. These conclusions concern sampling uncertainty;
they do not bound floating-point quadrature error or model error.
\end{proof}
